% Compiler: pdfLaTeX
\documentclass[11pt,a4paper]{article}

\usepackage[utf8]{inputenc}
\usepackage[T1]{fontenc}
\usepackage{lmodern}
\usepackage{microtype}
\usepackage{geometry}
\usepackage{enumitem}
\usepackage{titlesec}
\usepackage{fancyhdr}
\usepackage{hyperref}

% Document details
\newcommand{\researchtitle}{Network effects}
\newcommand{\researchshorttitle}{Network effects}
\newcommand{\researchdate}{2026-10-25}

% Layout
\geometry{left=2.7cm,right=2.7cm,top=2.5cm,bottom=2.5cm,headheight=14pt}
\hypersetup{hidelinks,pdftitle={\researchtitle}}
\urlstyle{same}
\frenchspacing
\setlength{\parindent}{0pt}
\setlength{\parskip}{0.55em}
\linespread{1.04}
\setlength{\emergencystretch}{1em}

\titleformat{\section}
    {\large\bfseries}{\thesection}{0.65em}{}
\titleformat{\subsection}
    {\normalsize\bfseries}{\thesubsection}{0.65em}{}
\titlespacing*{\section}{0pt}{1.15em}{0.35em}
\titlespacing*{\subsection}{0pt}{0.7em}{0.2em}
\setlist[itemize]{leftmargin=1.35em,itemsep=0.25em,topsep=0.2em}
\setlist[enumerate]{leftmargin=1.6em,itemsep=0.25em,topsep=0.2em}

\pagestyle{fancy}
\fancyhf{}
\fancyhead[L]{\small Random Research}
\fancyhead[R]{\small\researchshorttitle}
\fancyfoot[C]{\small\thepage}
\renewcommand{\headrulewidth}{0pt}
\renewcommand{\footrulewidth}{0pt}

\begin{document}

{\LARGE\bfseries\researchtitle\par}
\vspace{0.4em}
{\small\researchdate\par}

\section{Overview}

A product has network effects when its value to a user depends on how many
other people use it. A messaging app is an easy example: it becomes more
useful when the people you want to contact join. \cite{hbs}

That helps explain why choosing a service can feel like choosing a group.
The product's features are only part of what you get.

\section{Findings}

Direct effects operate within a user group. Another person joining a
communication service gives existing users another possible contact.
Indirect effects operate through another group or a complementary
product. More games can attract players to a console, while more players
give developers a reason to make games for it. \cite{hbs}

The distinction has consequences beyond terminology. Clements compares
models of direct and indirect effects and finds that their implications
for technological standardization can differ. His analysis is a reminder
to examine how the benefit arises before assuming that every network
behaves the same way. \cite{clements}

Positive feedback can make growth self-reinforcing, but it does not
guarantee a single winner. The HBS account also points to switching costs
and differences between platforms as factors in the outcome. \cite{hbs}

\section{Example}

Suppose two messaging apps are equally reliable. Most of a student's
class uses one, while the other has a feature the student prefers.
Choosing the first can be reasonable because the class is already there.

Now imagine a local marketplace with many buyers but almost no relevant
sellers. Its total user count may sound impressive without making it
useful for a particular purchase. In these hypothetical cases, I would
want to know who participates and what interaction becomes possible.

\section{Discussion}

I find network effects helpful for understanding why a service can be
hard to replace even when its interface is annoying. Moving would mean
getting other people to move as well.

What makes me skeptical is using the term as a substitute for explaining
a business. A growing user count leaves several questions unanswered.
Which users benefit from the newcomers? Can they actually interact?
Would the service remain useful if growth slowed?

For a small community, I would care more about having the right people
available than about membership across the whole platform.

\section{Practical implications}

When trying a service, I would ask what one more active user changes for
an existing user. If the answer is unclear, I would hold off on calling
its growth a network effect.

For a marketplace, I would examine the groups separately. A user count
alone tells me little about whether a buyer can find a suitable seller
when needed.

\renewcommand{\refname}{References}
\begin{thebibliography}{9}
\small
\setlength{\itemsep}{0.35em}

\bibitem{hbs}
Heinrich, A. (2026).
\href{https://online.hbs.edu/blog/post/what-are-network-effects}
{What are network effects?}
\textit{Harvard Business School Online}. October 22.

\bibitem{clements}
Clements, M. T. (2004).
\href{https://doi.org/10.1016/j.ijindorg.2004.01.003}
{Direct and indirect network effects: Are they equivalent?}
\textit{International Journal of Industrial Organization}, 22(5), 633--645.

\end{thebibliography}
\end{document}
